% TOP_PROBLEM: {"problemId": "problem.grothendieck-variational-hodge-conjecture", "canonicalTitle": "Grothendieck's Variational Hodge Conjecture", "releaseRank": 157, "primaryDomain": "geometry_topology", "primaryDomainLabel": "Geometry & topology", "displayStatus": "open_with_solved_subcases", "releaseStatus": "open_with_solved_subcases"}
% !TeX program = lualatex
% Definition and sources only; this file is not an LLM solution attempt.
\documentclass[11pt]{article}
\usepackage[margin=25mm]{geometry}
\usepackage{fontspec}
\setmainfont{Latin Modern Roman}
\usepackage{amsmath,amssymb,mathtools}
\usepackage{unicode-math}
\setmathfont{Latin Modern Math}
\usepackage{xurl,hyperref}
\hypersetup{colorlinks=true,urlcolor=blue}
\setcounter{secnumdepth}{0}
\setlength{\parindent}{0pt}
\setlength{\parskip}{5pt}
\setlength{\emergencystretch}{3em}
\begin{document}
\section*{\texorpdfstring{Grothendieck's Variational Hodge Conjecture}{Grothendieck's Variational Hodge Conjecture}}\label{157}
\subsection{Definitions and mathematical statement}
For a smooth projective family $f:X\to S$, the relative de Rham bundle
$\mathcal H^{2p}_{\mathrm{dR}}=\mathbf R^{2p}f_*\Omega^\bullet_{X/S}$ has the Gauss--Manin connection $\nabla$. A global algebraic section $\xi$ is horizontal when $\nabla\xi=0$, so analytically it is transported by the cohomological local system. Let $\mathrm{cl}_{\mathrm{dR}}$ be the cycle-class map from rational codimension-$p$ cycles on a fiber. The target is
\[
 \xi(s_0)\in\operatorname{im}\mathrm{cl}_{\mathrm{dR},s_0}
 \quad\Longrightarrow\quad
 \xi(s)\in\operatorname{im}\mathrm{cl}_{\mathrm{dR},s}\quad\text{for all }s\in S({\mathbb{C}}).
\]
Here $S$ is smooth and connected. No single relative cycle realizing all fibers is demanded; only fiberwise algebraicity of the transported class is asserted.
\par\smallskip\textit{Formulation and concept references:} \href{https://arxiv.org/pdf/1101.3647}{[S1]}.

\subsection{Short English statement}
If an algebraic cohomology class is transported horizontally through a smooth projective family, must it remain algebraic on every fiber?
\subsection{Sources}
\begin{itemize}
\item[S1]François Charles and Christian Schnell, Notes on absolute Hodge classes, Conjectures30/31, Proposition32, Theorem33 and Proposition34; base extension by the coordinator's recorded affine-open argument.
\url{https://arxiv.org/pdf/1101.3647}.
\textit{Formulation source.}
\end{itemize}
\end{document}
